\vspace{0.5em}\noindent\textbf{Create S3 bucket}\par

\begin{enumerate}
\def\labelenumi{\arabic{enumi}.}
\tightlist
\item
  Navigate to \textbf{S3 management console}
\item
  In the Bucket console, choose \textbf{Create bucket}
\end{enumerate}

\begin{figure}
\centering
\pandocbounded{\includegraphics[keepaspectratio,alt={Create bucket}]{5-Workshop/5.3-S3-vpc/create-bucket.png}}
\caption{Create bucket}
\end{figure}

\begin{enumerate}
\def\labelenumi{\arabic{enumi}.}
\setcounter{enumi}{2}
\tightlist
\item
  In \textbf{the Create bucket console}
\end{enumerate}

\begin{itemize}
\tightlist
\item
  \textbf{Name the bucket}: choose a name that hasn't been given to any
  bucket globally (hint: lab number and your name)
\end{itemize}

\begin{figure}
\centering
\pandocbounded{\includegraphics[keepaspectratio,alt={Bucket name}]{5-Workshop/5.3-S3-vpc/bucket-name.png}}
\caption{Bucket name}
\end{figure}

\begin{itemize}
\tightlist
\item
  Leave other fields as they are (default)
\item
  Scroll down and choose \textbf{Create bucket}
\end{itemize}

\begin{figure}
\centering
\pandocbounded{\includegraphics[keepaspectratio,alt={Create}]{5-Workshop/5.3-S3-vpc/create-button.png}}
\caption{Create}
\end{figure}

\begin{itemize}
\tightlist
\item
  Successfully create S3 bucket.
\end{itemize}

\begin{figure}
\centering
\pandocbounded{\includegraphics[keepaspectratio,alt={Success}]{5-Workshop/5.3-S3-vpc/bucket-success.png}}
\caption{Success}
\end{figure}

\vspace{0.5em}\noindent\textbf{Connect to EC2 with session
manager}\par

\begin{itemize}
\item
  For this workshop, you will use \textbf{AWS Session Manager} to access
  several \textbf{EC2 instances}. \textbf{Session Manager} is a fully
  managed \textbf{AWS Systems Manager} capability that allows you to
  manage your \textbf{Amazon EC2 instances} and on-premises virtual
  machines (VMs) through an interactive one-click browser-based shell.
  Session Manager provides secure and auditable instance management
  without the need to open inbound ports, maintain bastion hosts, or
  manage SSH keys.
\item
  First Cloud AI Journey
  \href{https://000058.awsstudygroup.com/1-introduce/}{Lab} for indepth
  understanding of Session manager.
\end{itemize}

\begin{enumerate}
\def\labelenumi{\arabic{enumi}.}
\tightlist
\item
  In the \textbf{AWS Management Console}, start typing
  \texttt{Systems\ Manager} in the quick search box and press
  \textbf{Enter}:
\end{enumerate}

\begin{figure}
\centering
\pandocbounded{\includegraphics[keepaspectratio,alt={system manager}]{5-Workshop/5.3-S3-vpc/sm.png}}
\caption{system manager}
\end{figure}

\begin{enumerate}
\def\labelenumi{\arabic{enumi}.}
\setcounter{enumi}{1}
\tightlist
\item
  From the \textbf{Systems Manager} menu, find \textbf{Node Management}
  in the left menu and click \textbf{Session Manager}:
\end{enumerate}

\begin{figure}
\centering
\pandocbounded{\includegraphics[keepaspectratio,alt={system manager}]{5-Workshop/5.3-S3-vpc/sm1.png}}
\caption{system manager}
\end{figure}

\begin{enumerate}
\def\labelenumi{\arabic{enumi}.}
\setcounter{enumi}{2}
\tightlist
\item
  Click \textbf{Start Session}, and select \textbf{the EC2 instance}
  named \textbf{Test-Gateway-Endpoint}.
\end{enumerate}

This EC2 instance is already running in ``VPC Cloud'' and will be used
to test connectivity to Amazon S3 through the Gateway endpoint you just
created (s3-gwe).

\begin{figure}
\centering
\pandocbounded{\includegraphics[keepaspectratio,alt={Start session}]{5-Workshop/5.3-S3-vpc/start-session.png}}
\caption{Start session}
\end{figure}

\textbf{Session Manager} will open a new browser tab with a shell
prompt: sh-4.2 \$

\begin{figure}
\centering
\pandocbounded{\includegraphics[keepaspectratio,alt={Success}]{5-Workshop/5.3-S3-vpc/start-session-success.png}}
\caption{Success}
\end{figure}

You have successfully start a session - connect to the EC2 instance in
VPC cloud. In the next step, we will create a S3 bucket and a file in
it.

\vspace{0.5em}\noindent\textbf{Create a file and upload to s3
bucket}\par

\begin{enumerate}
\def\labelenumi{\arabic{enumi}.}
\tightlist
\item
  Change to the ssm-user's home directory by typing
  \texttt{cd\ \textasciitilde{}} in the CLI
\end{enumerate}

\begin{figure}
\centering
\pandocbounded{\includegraphics[keepaspectratio,alt={Change user's dir}]{5-Workshop/5.3-S3-vpc/cli1.png}}
\caption{Change user's dir}
\end{figure}

\begin{enumerate}
\def\labelenumi{\arabic{enumi}.}
\setcounter{enumi}{1}
\tightlist
\item
  Create a new file to use for testing with the command
  \texttt{fallocate\ -l\ 1G\ testfile.\allowbreak{}xyz}, which will create a file of
  1GB size named ``testfile.xyz''.
\end{enumerate}

\begin{figure}
\centering
\pandocbounded{\includegraphics[keepaspectratio,alt={Create file}]{5-Workshop/5.3-S3-vpc/cli-file.png}}
\caption{Create file}
\end{figure}

\begin{enumerate}
\def\labelenumi{\arabic{enumi}.}
\setcounter{enumi}{2}
\tightlist
\item
  Upload file to S3 bucket with command
  \texttt{aws\ s3\ cp\ testfile.\allowbreak{}xyz\ s3:/\allowbreak{}/\allowbreak{}your-\allowbreak{}bucket-\allowbreak{}name}. Replace
  your-bucket-name with the name of S3 bucket that you created earlier.
\end{enumerate}

\begin{figure}
\centering
\pandocbounded{
\includegraphics[keepaspectratio,alt={Uploaded}]{5-Workshop/5.3-S3-vpc/uploaded.png}}
\caption{Uploaded}
\end{figure}

You have successfully uploaded the file to your S3 bucket. You can now
terminate the session.

\vspace{0.5em}\noindent\textbf{Check object in S3 bucket}\par

\begin{enumerate}
\def\labelenumi{\arabic{enumi}.}
\tightlist
\item
  Navigate to S3 console.\\
\item
  Click the name of your s3 bucket
\item
  In the Bucket console, you will see the file you have uploaded to your
  S3 bucket
\end{enumerate}

\begin{figure}
\centering
\pandocbounded{\includegraphics[keepaspectratio,alt={Check S3}]{5-Workshop/5.3-S3-vpc/check-s3-bucket.png}}
\caption{Check S3}
\end{figure}

\vspace{0.5em}\noindent\textbf{Section summary}\par

Congratulation on completing access to S3 from VPC. In this section, you
created a Gateway endpoint for Amazon S3, and used the AWS CLI to upload
an object. The upload worked because the Gateway endpoint allowed
communication to S3, without needing an Internet Gateway attached to
``VPC Cloud''. This demonstrates the functionality of the Gateway
endpoint as a secure path to S3 without traversing the Public Internet.
